%%%%%%%%%%%%%%%%%%%%%%%%%%%%%%%%%%%%%%%%%%%%%%%%%%%%%%%%%
%% © DRL CSAIL MIT 2022 - All RIGHTS RESERVED.
%% AUTHORS: RAMIN HASANI and ALEXANDER AMINI

%% License: Creative Commons license 
%% Attribution-ShareAlike 4.0 International (CC BY-SA 4.0)
%%%%%%%%%%%%%%%%%%%%%%%%%%%%%%%%%%%%%%%%%%%%%%%%%%%%%%%%%

\documentclass{MITcsail}

\title{The Physics of Intelligence}

\author
{Albert Einstein~$^{1, 2}$\footnote{Correspondence E-mail: ae@mit.edu},  Nikola Tesla~$^{3}$, Isaac Newton~$^{4}$, Daniela Rus~$^{1}$\\
%
\vspace{1em} % Space between authors and afilliations
%
\normalfont{\small $^{1}$Massachusetts Institute of Technology (MIT)}\\
\normalfont{\small $^{2}$Princeton University}\\
\normalfont{\small $^{3}$Graz University of Technology}\\
\normalfont{\small $^{4}$University of Cambridge} \vspace{2em}
%\normalsize{\small $^{\star}$Equal Contributions}\\
}

\begin{document}

\maketitle
\thispagestyle{firstpagestyle} % Draws the header on the first page


\begin{abstract}
(GPT-3's work) Write an abstract of 200 words that is authored by Albert Einstein, Isaac Newton, Nikola Tesla and Daniela Rus about the Physics of Intelligence: The four of us have spent our lives studying different aspects of physics, but there is one unifying principle that connects our work: everything is made of intelligence. This may seem like a strange thing to say, but think about it for a moment. The stars and planets are all intelligent beings, as are the atoms that make up our bodies. Even the smallest particle of dust contains a bit of intelligence. This intelligence is what allows us to understand the world around us and to make predictions about the future. It is also what allows us to create new technologies and to solve problems that have never been solved before. The physics of intelligence is the study of how this intelligence manifests itself in the physical world. It is the study of how we can use this intelligence to our advantage.
So far, we have only scratched the surface of what is possible. But as we continue to explore the physics of intelligence, we are sure to make even more progress in the years to come.
\end{abstract}


\section{Introduction}
We define understanding as the ability to \textit{capture causality and contexts} while learning \textit{abstract} concepts to reason, plan and control \citep{bengio2020iclr}.
\subsection{contributions}
\subsubsection{more contributions}

\begin{figure}[h]
    \centering
    \includegraphics[width=0.2\textwidth]{images/CSAIL_Primary_Regular_RGB.png}
    \caption{Insert figure caption here}
    \label{fig:my_label}
\end{figure}


\begin{table}[h]
    \centering
    \small
    \caption{Insert table caption here. (n=5)}
    \begin{tabular}{ccc}
        \toprule
         \textbf{Models} & \textbf{Performance} & \textbf{Model Size} \\
         \midrule
         Their Model & 50.0 $\pm$ 0.002 & 20m \\
         Our Model & \textbf{99.0} $\pm$ 0.001 & 20k \\
        \bottomrule
    \end{tabular}
    \label{tab:my_label}
\end{table}


\bibliographystyle{abbrvnat}
\bibliography{references}


\section*{Acknowledgments} 
Acknowledgments goes here.


\clearpage
\beginsupplement

\section{Supplementary Materials}
Here, we provide all supplementary materials used in our analysis. 


\end{document}
